% !TEX program = xelatex
% Short speaker script: 4 questions + pipeline + results
\documentclass[11pt,a4paper]{article}

\usepackage{fontspec}
\usepackage[hangul]{kotex}
\usepackage[margin=1.45cm]{geometry}
\usepackage{xcolor}
\usepackage{titlesec}
\usepackage{fancyhdr}
\usepackage{parskip}
\usepackage{amsmath,amssymb}
\usepackage[most]{tcolorbox}
\usepackage{paracol}

\setmainfont{Malgun Gothic}
\setsansfont{Malgun Gothic}

\definecolor{navy}{HTML}{1B365D}
\definecolor{teal}{HTML}{0F7B6C}
\definecolor{muted}{HTML}{5C6B7A}
\definecolor{coral}{HTML}{C45C3E}
\definecolor{sand}{HTML}{F7F5F0}

\pagestyle{fancy}
\fancyhf{}
\fancyhead[L]{\scriptsize\color{muted}FedSelfSelect script (KR / EN)}
\fancyhead[R]{\scriptsize\color{muted}short: problem / pipeline / results}
\fancyfoot[C]{\scriptsize\color{muted}\thepage}
\renewcommand{\headrulewidth}{0.3pt}

\titleformat{\section}
  {\large\bfseries\color{navy}}{\thesection}{0.7em}{}

\tcbuselibrary{breakable,skins}
\setlength{\columnsep}{16pt}
\setlength{\columnseprule}{0.45pt}
\colseprulecolor{navy!30}

\NewDocumentEnvironment{scriptslide}{mmm}{%
  \begin{tcolorbox}[
    colback=navy, colframe=navy, boxrule=0pt, arc=1pt,
    left=7pt,right=7pt,top=3.5pt,bottom=3.5pt,
    before skip=0.65em, after skip=0.25em,
    fontupper=\bfseries\color{white}\normalsize,
  ]
  Slide #1.\quad #2\hfill{\mdseries\scriptsize #3}
  \end{tcolorbox}
  \nopagebreak
  \begin{paracol}{2}
  \small
  \parskip=0.38em\relax
  \emergencystretch=2em\relax
  {\color{navy}\textbf{한글.}}\par
}{\end{paracol}}

\newcommand{\key}[1]{\par\vspace{0.22em}{\color{coral}\textbf{#1}}\par}

\NewDocumentCommand{\en}{+m}{%
  \switchcolumn
  {\color{teal}\textbf{English.}}\par
  #1%
}

\begin{document}

\begin{center}
{\LARGE\bfseries\color{navy}FedSelfSelect 발표 대본}\\[0.3em]
{\large\color{teal}Identification-free rare-class learning}\\[0.35em]
{\footnotesize\color{muted}읽을 파일: \texttt{FedSelfSelect.pdf} \quad 영어는 오른쪽}
\end{center}

\vspace{0.2em}
\noindent
발표는 \textbf{오른쪽 English}를 읽습니다. 왼쪽은 같은 내용입니다.
주황 문장만 말해도 흐름은 됩니다. 전체 약 7분입니다.

\begin{tcolorbox}[colback=sand,colframe=teal,boxrule=0.5pt,arc=1pt,
  left=7pt,right=7pt,top=5pt,bottom=5pt]
{\color{navy}\textbf{이 네 질문.}}
1. Problem $\rightarrow$ 슬라이드 1.\\
2. Equations $\rightarrow$ 슬라이드 3.\\
3. Experiment setup $\rightarrow$ 슬라이드 5.\\
4. Loss / theory $\rightarrow$ 슬라이드 4.\\
파이프라인은 슬라이드 2. 결과는 슬라이드 6.
\end{tcolorbox}

\section{Problem}

\begin{scriptslide}{1}{What issue in FL?}{약 70초}
질문 1입니다. 문제 정의입니다.

연합분류는 데이터 $K$명에게 두고 공유 모델 $w$ 하나를 학습합니다.
Stratify 같은 라벨 인식 방법은 클래스마다 턴을 줍니다.
서버가 클래스 $c$ 홀더를 지명해야 그 턴이 돕니다.

희귀 클래스는 홀더가 2명뿐입니다.
라벨 메타데이터에 노이즈를 넣고 그걸로 사람을 고르면 희귀가 먼저 깨집니다.

pp-CFL 비트는 없는 사람을 있다로 만듭니다. Covertype 희귀 1.2\%, overall은 62\%라 안 보입니다.
Stratify 카운트는 홀더 30장을 살리고 희귀 70.7\%입니다. 대신 서버가 그 2명을 압니다.
FedAvg는 지명을 안 하지만 라벨도 안 봐서 희귀가 0입니다.

\key{문제는 Stratify보다 희귀를 올리는 것이 아닙니다. 서버가 홀더를 고르지 않고 희귀를 학습하는 것입니다.}
\en{Question 1. The problem.

Federated classification keeps data on $K$ clients and fits one shared model $w$.
Label-aware methods such as Stratify give each class a turn.
The server must name a holder of class $c$, or that turn does not run.

The rare class has only two holders.
If the server adds noise to label metadata and then routes from it, the rare class breaks first.

pp-CFL bits turn non-holders into false holders. Covertype rare is 1.2 percent; overall stays 62 percent, so the failure is hidden.
Stratify counts keep the thirty samples and rare reaches 70.7 percent, but the server named those two clients.
FedAvg names nobody and inspects no labels, so rare is zero.

\key{The issue is not to beat Stratify on rare accuracy. It is to learn the rare class without the server choosing holders.}}
\end{scriptslide}

\section{Algorithm}

\begin{scriptslide}{2}{Pipeline, nine steps}{약 80초}
알고리즘입니다. 1--2는 한 번, 3--8은 한 스텝, 9는 1500번입니다.

1. 각 클라이언트가 클래스마다 있다/없다만 올립니다. 서버는 합 $n_c$만 압니다. 누구인지는 모릅니다.

2. $n_c$에 라플라스 노이즈를 붙여 $\tilde n_c$를 만듭니다. 노이즈는 이 숫자에만 붙습니다.

3. 서버가 모델 $w$를 뿌리고 ``이번 클래스는 $c$''라고만 말합니다. 사람 이름을 고르지 않습니다.

4. 클래스 $c$가 있는 클라이언트는 그 샘플로 그래디언트를 계산해서 항상 올립니다.

5. 없는 클라이언트는 확률 $q$로만 0을 올립니다. 나머지는 침묵합니다.

6. 서버는 SecAgg 합만 받습니다. 값이 그래디언트인지 0인지, 누가 올렸는지는 모릅니다.

7. 그 합을 올린 사람 수가 아니라 홀더 수 $\tilde n_c$로 나눕니다. 0은 합을 바꾸지 않습니다.

8. 클래스 $1,\ldots,C$를 같은 방식으로 돈 뒤, 그 평균으로 모델을 한 번 업데이트합니다.

9. 3--8을 반복합니다. 1--2는 다시 하지 않습니다.

\key{서버는 홀더를 고르지 않습니다. 합만 받고, 홀더 수로 나눕니다.}
\en{The algorithm. Steps 1--2 once. Steps 3--8 one step. Step 9 repeats that 1500 times.

1. Each client sends only yes or no per class. The server sees the sum $n_c$, not who.

2. Laplace noise on $n_c$ gives $\tilde n_c$. Noise sits only on that number.

3. The server broadcasts $w$ and says only ``this class is $c$.'' It does not pick a name.

4. A client that has class $c$ computes a gradient on those samples and always uploads.

5. A client that does not have class $c$ uploads 0 with probability $q$. The rest stay silent.

6. The server receives only the SecAgg sum. It does not see zeros versus gradients, and it does not see who sent them.

7. It divides by the holder count $\tilde n_c$, not by the number of uploaders. Zeros do not change the sum.

8. After all $C$ classes, it updates the model once with that average.

9. Repeat 3--8. Do not repeat 1--2.

\key{The server does not choose holders. It sees a sum and divides by the holder count.}}
\end{scriptslide}

\section{Equations}

\begin{scriptslide}{3}{Key equations}{약 70초}
질문 2입니다. 알고리즘 말고 식입니다.

(1) $\mathcal{H}_c$는 홀더 집합, $n_c$는 인원. 서버가 봐도 되는 것은 합입니다.

(2) $\tilde n_c = \max(n_c+\mathrm{Lap}(1/\varepsilon),1)$. 나눗셈 분모입니다. 라우팅 점수가 아닙니다.

(3) 홀더 그래디언트 $g_k=\nabla\ell(w;\mathcal{B}_k)$.

(4) 없는 사람은 확률 $q$로만 0을 올립니다.

(5) $S_c$는 SecAgg 합입니다.

(6) $\bar g_c=S_c/\tilde n_c$. $|\mathcal{U}_c|$로 나누면 안 됩니다.

(7) $w\leftarrow w-\eta(\sum_c\bar g_c)/C$. 클래스 균등 평균입니다. $n_k$ 가중이 아닙니다.

(8) 업로드 $\mathbb{E}[|\mathcal{U}_c|]=n_c+q(K-n_c)$. 희귀 턴은 약 12명입니다.

\key{핵심은 식 (6)입니다. 홀더 수로 나눕니다. 올린 사람 수로 나누지 않습니다.}
\en{Question 2. The equations besides the algorithm.

(1) $\mathcal{H}_c$ is the holder set, $n_c$ the count. The server may see a sum, not a name list.

(2) $\tilde n_c=\max(n_c+\mathrm{Lap}(1/\varepsilon),1)$. A divisor, not a routing score.

(3) Holder gradient $g_k=\nabla\ell(w;\mathcal{B}_k)$.

(4) A non-holder uploads 0 with probability $q$.

(5) $S_c$ is the SecAgg sum.

(6) $\bar g_c=S_c/\tilde n_c$. Do not divide by $|\mathcal{U}_c|$.

(7) $w\leftarrow w-\eta(\sum_c\bar g_c)/C$. Unweighted mean over classes, not an $n_k$ average over clients.

(8) Uploads $\mathbb{E}[|\mathcal{U}_c|]=n_c+q(K-n_c)$. About twelve on a rare turn.

\key{Equation (6) is the key. Divide by the holder count, not by the number of uploaders.}}
\end{scriptslide}

\section{Loss and theory}

\begin{scriptslide}{4}{Loss and guarantee}{약 50초}
질문 4입니다.

손실은 softmax 교차엔트로피입니다.
$\ell(w;x,y)=-\log[\mathrm{softmax}(f_w(x))]_y$.
RMSE가 아닙니다. 모델은 hidden 64 ReLU MLP입니다.
점수는 overall, 희귀 recall, balanced accuracy입니다.

이론은 항등식 두 개입니다.
Prop.~1. dummy 0은 합을 안 바꿉니다. $S_c$는 홀더 합입니다.
Prop.~2. 라플라스는 분모만 흔듭니다. 누구를 학습시킬지는 안 고릅니다.
$n_c$ 카운트만 $\varepsilon$-DP입니다. $w_T$의 DP도, 수렴도, 희귀 accuracy 하한도 없습니다.
통신이 큰 것은 limitation입니다.

\key{손실은 CE입니다. 보장은 합이 홀더 합이라는 것입니다. 희귀 점수 정리는 없습니다.}
\en{Question 4.

The loss is softmax cross-entropy: minus log of the true-class softmax probability.
Not RMSE. The model is a ReLU MLP of width 64.
We report overall accuracy, rare recall, and balanced accuracy.

Theory is two identities.
Prop.~1. Dummy zeros leave $S_c$ equal to the holder sum.
Prop.~2. Laplace perturbs only the divisor. It does not choose who trains.
Only the count $n_c$ is $\varepsilon$-DP. We do not claim DP of $w_T$, convergence, or a rare-accuracy bound.
Large communication is a limitation.

\key{The loss is cross-entropy. The guarantee is that the sum is the holder sum. There is no rare-accuracy theorem.}}
\end{scriptslide}

\section{Setup}

\begin{scriptslide}{5}{Experiment setup}{약 55초}
질문 3입니다.

Covertype, Shuttle, Dry Bean. 표 분류입니다.
$K=100$. 홀더는 그 클래스 샘플이 있는 클라이언트입니다.
희귀 홀더 2명, 각자 희귀 샘플 30장. 30은 배치가 아닙니다. 그 2명의 희귀 데이터 크기입니다.
흔한 클래스는 클라이언트당 2--3개, 미니배치 64. 희귀는 30장 전부 씁니다.
$\varepsilon=0.5$, $q=0.1$, seed 3개. 이 나눔은 실험 세팅입니다.

시계는 두 개입니다. 1500과 $150\times 10\times 10$은 같은 값이 아닙니다.
FedSelfSelect, Stratify, pp-CFL, Oracle은 클래스 턴 1500스텝. 공정 비교는 이 네 개입니다.
FedAvg, FedProx, SCAFFOLD, FedDyn, MOON은 150 라운드 $\times$ 10명 $\times$ 10 로컬 스텝. 라벨을 안 보는 참조입니다.

\key{홀더 2명, 30장은 데이터 크기입니다. 공정 비교는 클래스 턴 1500끼리입니다.}
\en{Question 3.

Covertype, Shuttle, Dry Bean. Tabular classification.
$K=100$. A holder is a client who has that class.
Two rare holders, thirty rare samples each. Thirty is not the batch size. It is their rare dataset size.
Each client holds two or three common classes. Common minibatches are 64. Rare uses all thirty.
$\varepsilon=0.5$, $q=0.1$, three seeds. This split is an experimental setup.

Two clocks. 1500 is not $150\times 10\times 10$.
FedSelfSelect, Stratify, pp-CFL, and Oracle share 1500 class-turn steps. That comparison is fair.
FedAvg, FedProx, SCAFFOLD, FedDyn, and MOON share 150 rounds $\times$ 10 clients $\times$ 10 local steps. They are label-blind references.

\key{Two holders. Thirty is dataset size. The fair comparison is the 1500 class-turn methods.}}
\end{scriptslide}

\section{Results}

\begin{scriptslide}{6}{Accuracy}{약 55초}
Covertype이 증거가 갈리는 표입니다. overall은 다 약 58--62\%입니다.

희귀만 갈립니다.
FedAvg, FedProx, SCAFFOLD, FedDyn, MOON은 전부 0.0입니다. 라벨을 몰라서 홀더 2명을 거의 안 봅니다.
pp-CFL 1.2. 비트가 홀더를 지웁니다.
Stratify 70.7. 카운트는 살지만 홀더를 지명합니다.
FedSelfSelect 77.3. 지명 없이 희귀가 들어옵니다.
Oracle 79.6이 상한입니다.

70.7 대 77.3이 이유가 아닙니다. 이유는 식별을 안 한다는 것입니다.
Shuttle도 희귀가 갈립니다. Dry Bean은 BOMBAY가 쉬워서 방법이 거의 안 갈립니다.
통신은 limitation입니다. $q=0.1$이어도 Covertype에서 Stratify보다 약 42배입니다.

\key{overall은 숨깁니다. Covertype 희귀가 문제입니다. 주장은 식별 없이 희귀를 학습하는 것입니다.}
\en{Covertype is the table that splits the claim. Overall is about 58 to 62 percent for every method.

Only rare moves.
FedAvg, FedProx, SCAFFOLD, FedDyn, and MOON are all 0.0. They never inspect labels, so they almost never see the two holders.
pp-CFL is 1.2. Bits erase the holders.
Stratify is 70.7. Counts survive, but the server named the holders.
FedSelfSelect is 77.3. The rare class enters without naming anyone.
Oracle 79.6 is the ceiling.

The 70.7 versus 77.3 gap is not why we use the method. The reason is no identification.
Shuttle still splits on rare. Dry Bean does not, because BOMBAY is easy.
Communication is a limitation. Even at $q=0.1$, Covertype uploads are about 42 times Stratify.

\key{Overall hides the failure. Covertype rare is the object. The claim is rare learning without naming holders.}}
\end{scriptslide}

\end{document}
